\documentclass[a4paper,12pt]{article}

\usepackage{cmap}		
\usepackage[utf8]{inputenc}			
\usepackage[english,russian]{babel}
\usepackage{framed}
\usepackage{hyperref}
\usepackage{amsmath}
\usepackage{graphicx}
\usepackage[colorinlistoftodos]{todonotes}
\usepackage{wrapfig}
\usepackage{lipsum}
\usepackage{listings}
\usepackage{color}
\usepackage{indentfirst}
%\usepackage{times}
\usepackage{textcomp}
\usepackage{indentfirst}
\textheight=24cm % высота текста
\textwidth=16.5cm % ширина текста
\oddsidemargin=0pt % отступ от левого края
\topmargin=-1.5cm % отступ от верхнего края
\definecolor{mygray}{rgb}{0.4,0.4,0.4}
\definecolor{mygreen}{rgb}{0,0.8,0.6}
\definecolor{myorange}{rgb}{1.0,0.4,0}

\lstdefinestyle{customc}{
  belowcaptionskip=1\baselineskip,
  breaklines=true,
  frame=L,
  xleftmargin=\parindent,
  language=C,
  showstringspaces=false,
  basicstyle=\footnotesize\ttfamily,
  keywordstyle=\bfseries\color{green!40!black},
  commentstyle=\itshape\color{purple!40!black},
  identifierstyle=\color{blue},
  stringstyle=\color{orange},
  numbers=left,
  numbersep=13pt,
  numberstyle=\small\color{mygray},
}
\lstset{escapechar=@,style=customc}
\linespread{1.5}
\newcommand{\HRule}{\rule{\linewidth}{0.5mm}}

\begin{document}

\begin{titlepage}
\begin{center}

\textsc{\Large Московский Физико-Технический Институт}\\
\textsc{\large (Национальный Исследовательский Университет)}\\[1.5cm]

% Upper part of the page. The '~' is needed because \\
% only works if a paragraph has started.
\includegraphics[width=0.25\textwidth]{img/logo.png}~\\[1cm]

\textsc{\Large Бакалаврская работа}\\[0.5cm]

% Title
\HRule \\[0.4cm]
{ \LARGE \bfseries Влияние вышестоящих регуляторных регионов и промоторов на связь между бактериофагами и их хозяевами \\[0.4cm] }

\HRule \\[1.5cm]

% Author and supervisor
\noindent
\begin{minipage}{0.4\textwidth}
\begin{flushleft} \large
\emph{Студент:}\\
Шарафутдинов Эмиль
\end{flushleft}
\end{minipage}%
\begin{minipage}{0.4\textwidth}
\begin{flushright} \large
\emph{Преподаватель:} \\
 Родригес Валера Франсиско Эдуардо
\end{flushright}
\end{minipage}

\vfill

% Bottom of the page
{\large 24 апреля 2021 г.}

\end{center}
\end{titlepage}

\newpage
\begin{center}
\item \section{Аннотация}
\end{center}

\par{Бактериофаги (фаги) - это вирусы, которые избирательно поражают бактериальные клетки. Такая специфика обусловлена
зависимостью фага от клеточного механизма хозяина, который нужен фагу для репликации и биосинтеза вирусных компонентов.
Пример такой зависимости  можно привести для фага Т4. Во время транскрипции, ранняя конкуренция за РНК-полимеразу
хозяина позволяет вирусу полностью захватить механизм репликации путем сигма-присвоения(sigma appropriation) [1].
Недавно показали[2], что фаги, которые заражают цианобактерии, кодируют ген фотосинтеза. Причём транскрипция этого гена
реагирует на те же условия окружающей среды, что и ген, который кодируется у хозяина. Из этого можно предположить, что
существует связь между регуляторными элементами фага и регуляторными элементами их хозяина. Это поможет лучше
предсказывать взаимоотношения фаг-хозяин.}

\newpage
\tableofcontents



\newpage
\begin{center}
\item \section{Введение}
\end{center}

    \par{В одной статье были собраны результаты разных исследований, касающихся взаимоотношений фага T4 и его хоста – E.
    Coli. В этих работах было показано, как ведёт себя фаг, который внедряется в бактерию.
	У фага Т4 нет собственных РНК-полимераз, поэтому он использует РНК-полимеразы хоста для транскрипции. Но почему РНК
	полимераза хоста должна сразу транскрибировать именно с ДНК Т4? Не должна. Но всё же обнаружили, что некоторая часть
	ДНК фага транскрибируется РНК-полимеразой хоста сразу. Эту часть транскрипции назвали ранней транскрипцией фага. На
	этом слайде можно видеть, почему эта ранняя транскрипция просходит. Здесь обозначены разные мотивы промоторов ДНК
	фага (посередине – ранняя, снизу - промежуточная) и хоста (сверху). Эти мотивы, которые находятся в промоторе,
	называются вышестоящими регуляторными элементами. Как видно, в промоторе ранней транскрипции регуляторные элементы
	между фагом и хостом совпадают, и выше -35 участка, что позволяет РНК полимеразе цепляться за ДНК фага также хорошо,
	как за ДНК хоста.  Всё, что до мотива TGn совпадает у всех промоторов. MotA Box в промоторе промежуточной
	транскрипции нужен для зацепки белка, который способствует промежуточной транскрипции. Вкратце, что там просходит: в
	этот момент с РНК, транскрибировавшейся от фага (в ранней транскрипции), транслируются белки, которые модифицируют
	РНК-полимеразы хозяина таким образом, что далее начинают транскрибироваться только ДНК фага. В том числе в начальной
	транскрипции ещё образуется белок MotA, который позволяет ДНК прикрепиться к модифицированной РНК-полимеразе. После
	чего с этой ДНК начинает транскрибироваться РНК. Этот этап уже называют промежуточной транскрипцией. Этот процесс
	присвоения РНК-полимеразы хозяина называют сигма присвоением. Ещё раз. Проиcходит это так:}

	\par{В самом начале транслируются белки, которые глушат РНК-полимеразу Вместе с этим транслируются белки, которые
	позволяют при заглушенной полимиразе транскрибироваться ДНК фага на этой полимеразе, они садятся на ДНК фага. Таким
	образом фаг полностью захватывает РНК-полимеразу хозяина Рассмотри ещё один интересный случай для нас. Ещё в одной
	статье авторы смотрели, как ведут себя транкриптомы фагов и их хозяев (цианобактерий в этом случае) в зависимости от
	интенстивности света, которым они светили. При маленькой интенсивности транскриптомы делились в среднем поровну, но
	при бОльшей интенсивности среднее число РНК фагов было намного больше, чем у цианобактерий. Авторы решили
	посмотреть, что это за участок РНК, который так сильно возрастает у фага при сильном свете. Они обнаружили, что это
	участок РНК, ответственный за белок psbA, который оказался общим у фага и его хозяина. Авторы предположили, что это
	вполне возможно, если фаги позаимствовали гены у хозяев в ходе рекомбинации.}
	
    \begin{center}
    \item \subsection{English Wiki}
    \end{center}

	\par{Бактериофаг , также известный неофициально как фаг, это вирус, который
	заражает и размножается внутри бактерии и археи. Термин произошел от "бактерии" и греческого (фагеин), что
	означает "пожирать". Бактериофаги состоят из белков, которые инкапсулируют геном ДНК или РНК, и могут иметь
	структуры, которые либо просты, либо сложны. Их геномы могут кодировать всего четыре гена (например, MS2) и сотни
	генов. Фаги размножаются внутри бактерии после введения своего генома в ее цитоплазму.}
	
	\par{Бактериофаги являются одними из наиболее распространенных и разнообразных организмов в биосфере.[1]
	Бактериофаги-это вездесущие вирусы, встречающиеся везде, где существуют бактерии. Подсчитано, что на планете
	насчитывается более 1031 бактериофага, больше, чем любой другой организм на Земле, включая бактерии, вместе
	взятые.[2] Вирусы являются наиболее распространенным биологическим существом в водной толще мирового океана и вторым
	по величине компонентом биомассы после прокариот, [3] где в микробных матах на поверхности было обнаружено до 9x108
	вирионов на миллилитр, [4] и до 70\% морских бактерий могут быть заражены фагами.[5] Англоязычная вики}

    \par{Фаги используются с конца 20 - го века в качестве альтернативы антибиотикам в бывшем Советском Союзе и
    Центральной Европе, а также во Франции.[6][7] Они рассматриваются как возможная терапия против штаммов многих
    бактерий с множественной лекарственной устойчивостью (см. фаготерапию).[8] С другой стороны, было показано, что фаги
    Inoviridae усложняют биопленки, участвующие в пневмонии и муковисцидозе, и защищают бактерии от лекарств,
    предназначенных для искоренения болезни, тем самым способствуя персистирующей инфекции[9].]}
    
\begin{center}
   \item \subsection{Bacteriophage genomics \cite{genomics}}
\end{center}

    
\begin{center}
   \item \subsection{Statistical structure of host–phage interactions \cite{structure}}
\end{center}
    

   
\begin{center}
   \item \subsection{Phage–bacteria infection networks \cite{networks}}
\end{center}

    
\newpage
\begin{center}
\item \section{Цель и задачи} \label{sec:code}
\end{center}

\par{Эти статьи – хорошие примеры, но они не единственные в своём роде. Нам стали интересны такие события, и в нашей
    работе мы предположили, что такие схожие вышестоящие регуляторные участки в ДНК фагов и их хостов могут встречаться
    не только у цианобактерий и их фагов и T4 с E. Coli, а у всех фагов и их хостов. В том числе, мы посчитали, что у
    многих фагов и их хозяев могут встречаться общие белки. Таким образом, если мы сможем установить взаимоотношения
    фагов и их хозяинов, мы в будущем сможем прогнозировать отношения фагов и каких-либо хостов.}

\newpage
\begin{center}
\item \section{Обзор литературы} \label{sec:math}
\item \subsection{Hinton D.M. (Обработанный)}
\end{center}
        
        \par{Экспрессия генома Т4 - это строго регулируемый процесс, который начинается сразу после заражения хозяина. 
        Основной контроль над этой экспрессией происходит на уровне транскрипции. Т4 не кодирует свою собственную 
        РНК-полимеразу (RNAP), а вместо этого кодирует множество факторов, которые служат для изменения специфичности 
        полимеразы хозяина по мере развития инфекции. Изменения полимеразы хозяина сопровождаются с тремя классами 
        транскрипции фага: ранней, промежуточной и поздней. Ранняя и промежуточная РНК (РНК, образовавшаяся при ранней и
        промежуточной транскрипции) обнаруживается дорепликативно 
        \cite{hinton1,hinton2,hinton3,hinton4,hinton5,hinton6}, в то время как поздняя транскрипция сопровождается 
        репликацией Т4 и нам в нашей работе не интересна. Ранние транскрипты T4 генерируются из ранних промоторов 
        (Pe),которые активны сразу после инфицирования. Промежуточные транскрипты Т4 генерируются примерно через 1 
        минуту после заражения при \(37^\circ C\) и требуют синтеза белка фага. Средняя РНК синтезируется двумя 
        способами: 1) активация промежуточных промоторов (Pm) и 2) расширение транскриптов Pe из ранних генов в 
        нижестоящие средние гены.}
        
        \begin{center}
          \item {\Large Механизм начала транскрипции кишечной палочки}
        \end{center}
        
        \par{E. coli RNAP холофермент, как и все бактериальные RNAP, состоит из ядра субъединиц 
        \((\beta,\beta',\alpha_1,\alpha_2\) и \(\omega)\), которое содержит активный сайт для синтеза РНК, и фактора 
        специфичности \(\sigma\), который распознает промоторы в ДНК и устанавливается на начальный сайт для 
        транскрипции. Первичный \(\sigma\), \(\sigma^{70}\) в E. coli, используется во время экспоненциального роста; 
        альтернативные факторы \(\sigma\) направляют транскрипцию генов, необходимых в различных условиях роста или во 
        время стресса [рассмотрено в [17-19]]}
        
        \par{Для начала транскрипции части RNAP должны сначала распознать и связаться с двухцепочечными (ds) элементами 
        распознавания ДНК, присутствующими в промоторе [reviewed in [20]]. Каждый из C-концевых доменов 
        $\alpha$-субъединиц ($\alpha$-CTDS) может взаимодействовать с вышестоящим элементом, A/T богатыми 
        последовательностями, присутствующими между позициями -40 и -60. Части \(\sigma^{70}\), присутствующие в RNAP, 
        могут взаимодействовать с тремя различными элементами dsDNA: с элементом -35, с последовательностью -15TGn-13 
        (TGn) и с позициями -12/-11 элемента -10. Как правило, промотор должен содержать только два из трех 
        \(\sigma^{70}\)-зависимых элементов для активности; таким образом, промоторы E. coli могут быть свободно 
        классифицированы как -35/-10 (основной класс), TGn/-10 (также называемый расширенным -10) или -35/TGn 
        [рассмотрено в [20]].}
        
        \begin{center}
            \item {\Large Ранняя транскрипция Т4}
        \end{center}
        
        \par{Т4 заражает E.coli только во время экспоненциального роста. Транскрипция ранних генов Т4 начинается сразу 
        после заражения. Таким образом, для эффективной инфекции фаг должен быстро перенаправить 
        \(\sigma^{70}\)-ассоциированный RNAP, который активно участвует в транскрипции генома хозяина, на ранние 
        промоторы T4. Такое немедленное перенаправление часто происходит успешно отчасти потому, что большинство ранних 
        промоторов T4 содержат такие же элементы распознавания \(\sigma^{70}\)-RNAP (-35, TGn и -10 элементов) и 
        элементы \(\alpha\)-CTD UP (списки известных последовательностей ранних промоторов T4 см. в [4,5]). Тем не менее, выравнивание последовательностей ранних промоторов Т4 показало дополнительные области совпадения, и есть предположение, что они содержат и другие последовательности, которые могут оптимизировать взаимодействие RNAP хозяина с элементами промотора T4. Следовательно, в отличие от большинства промоторов хозяина, которые относятся к классам -35/-10, TGn/-10 или -35/TGn, ранние промоторы T4 могут быть описаны как "супер" UP/-35/TGn/-10 промоторы \cite{hinton}. Действительно, большинство ранних промоторов Т4 очень хорошо конкурируют с промоторами хозяина за доступный RNAP [39].}
        
        \begin{center}
            \item {\Large Продолжение ранних стратегий доминирования Т4}
            \par{Здесь в основном про синтез белков - подумать потом, надо ли это включать}
        \end{center}
        
        \begin{center}
            \item {\Large Транскрипция средних генов путем расширения ранних транскриптов}
            \par{Надо ли это вообще включать?}
        \end{center}
        
        
\begin{center}
    \item \subsection{Energy limitation of cyanophage development: implications formarine carbon cycling}
\end{center}
       
\newpage
\begin{center}
\item \section{Материалы и Методы} \label{sec:code}
\item \subsection{Программные средства}
\end{center}
\begin{itemize}
    \item Virus-Host DB - База данных Virus-Host организует данные о взаимоотношениях между вирусами и их хостами,
    представленные в виде пар идентификаторов таксономии NCBI для вирусов и их хостов. База данных Virus-Host охватывает
    вирусы с полными геномами, хранящимися в 1) NCBI/RefSeq и 2) GenBank, номера присоединения которых указаны в геномах
    EBI. Информация о хосте собирается из RefSeq, GenBank (в свободном текстовом формате), UniProt, ViralZone и вручную
    курируется с дополнительной информацией, полученной в результате опросов литературы. \cite{virus-host}
    
    \item GeneMarkS-2 - программное обеспечение для поиска генов внутри геномов. \cite{lomsad}
    
    \item NCBI Entrez Programming Utilities - программное обеспечение, набор из восьми серверных программ, которые
    обеспечивают стабильный интерфейс с системой запросов и баз данных Entrez в Национальном центре биотехнологической
    информации (NCBI) \cite{entrez}
    
    \item Biopython - набор свободно доступных инструментов для биологических вычислений, написанных на Python
    международной командой разработчиков. \cite{biopython}
    
    \item The MEME Suite - программное обеспечение для поиска новых, нераскрытых мотивов (повторяющихся паттернов
    фиксированной длины) в последовательностях. МЕМ разбивает паттерны переменной длины на два или более отдельных
    мотива. \cite{bailey}
    
    \item Python 3.8.5 - среда для запуска скриптов формата *.py \cite{python}
    
    \item Jupyter Notebook -веб-приложение с открытым исходным кодом, которое позволяет создавать и обмениваться
    документами, содержащими живой код, уравнения, визуализации и повествовательный текст. Использование включает в себя:
    очистку и преобразование данных, численное моделирование, статистическое моделирование, визуализацию данных, машинное
    обучение и многое другое. \cite{jupyter}
    
    \item Скрипты отсюда \(https://github.com/poolsar42/phages-and-hosts\) - подумать, что про них написать!
    \cite{github}
    
    \item \(Fasta\_merge.ipynb\) - скрипт для объединения регуляторных областей хозяина и его фагов в один файл.
    \item \(Hosts\_upstream\_regions.ipynb\) - скрипт для выделения регуляторных областей.
    \item \(Proteins\_formation.ipynb\) - скрипт для сбора белковых последовтельностей в один файл.
    \item \(Searching\_host\_genomes.ipynb\) - скрипт для выгрузки необходимых к работе геномов хостов.
    \item \(Unpacking.ipynb\) - скрипт для распаковки выгруженных геномов.
    \item \(VH\_tsv\_writer.ipynb\) - скрипт, выписывающий в файл ID хоста и ID прилежащих к нему фагов.
    \item \(Virus\_genes.ipynb\) - скрипт для выгрузки геномов фагов.
    \item \(missing\_phage\_genomes.ipynb\) - скрипт для поиска фагов, для которых GMS2 не находит гены.
    \item \(proteins\_convertation.ipynb \) - скрипт для преобразования \\ ДНК-последовательностей
                                в белковые.
    \item \(regex\_promoter.ipynb \) - скрипт для анализа полученных результатов
    \item \(run\_gms2.py\) - скрипт для поиска генов в геномах
    \item \(tsv\_write.ipynb\) - скрипт для выписывания информации о связи между хостами и фагами в один файл
\end{itemize}

\begin{center}
\     item \subsection{Получение общих мотивов}
      \item  \subsubsection{Протокол выгрузки геномов}
\end{center}
    \begin{itemize}
        \item С FTP сервера Virus-Hist DataBase выгрузить файл virushostdb.tsv - в нём содержится информация о связи между фагами и их хостами.
    
        \item Запустить скрипт \(tsv\_write.ipynb\) - выписывает в удобном формате все взаимоотношения между фагами и хостами.
    
      \item Запустить скрипт \(Virus\_genes.ipynb\) - для выгрузки геномов фагов.
        \item Запустить скрипт \(Searching\_host\_genomes.ipynb\) - для выгразки геномов хостов.
        \item Запустить скрипт \(Unpacking.ipynb\).

    \end{itemize}

    \begin{center}
    \item \subsubsection{Протокол получение вышестоящих регуляторных регионов}
    \end{center}
        \begin{itemize}
            \item Запустить скрипт \(run\_gms2.py\)
            \item Запустить скрипт \(Hosts\_upstream\_regions.ipynb\)
        \end{itemize}
    \begin{center}
    \item \subsubsection{Протокол подготовки данных к поиску мотивов}
    \end{center}
        \begin{itemize}
            \item Запустить скрипт \(Fasta\_merge.ipynb\)
            \item Запустить MEME Suite с получившимися файлами.
        \end{itemize}

\par{Нам нужно было посмотреть на сходтсва между вышестоящими регуляторными участками у хостов и фагов. Для этого нужно
было посмотреть на геномы хостов и фагов. Мы собрали геномы всех фагов, которые находятся в базе данных Virus-Host DB,
там для каждого фага также был указан или были указаны хосты, к которым он может прикрепиться. Их RefSeq ID. Мы всё
выгрузили с FTP сервера.}

\par{Далее мы выгрузили все геномы хостов с базы данных RefSeq NCBI. Когда у нас были геномы фагов и их хостов, нам было
нужно получить с этих геномов гены. Гены мы находили с помощью Gene-Marks 2. Что хорошо, Gene-Marks 2 также дал позицию
каждого гена в геноме. Каждому гену посчитали его транслирующийся белок, также с помощью Gene-Marks 2. После этого для
каждого гена нашли вышестоящие регуляторные области, длиной в 50 нуклеотидов, которые нужны для посадки РНК-полимераз(?).
Наша гипотеза как раз заключается в том, что эти вышестоящие области у фага и его хозяина для некоторых генов будут
совпадать по некоторым мотивам.}

\par{Мы поработали с последовательностями. Мы выгружали геномы всех фагов и хостов отдельными файлами. Мы предварительно
обработали файлы, достали из них вышестоящие регуляторные области длинной в 50 нуклеотов. Вышестоящие области хоста и
всех фагов, прилегающих к нему, мы закинули в один файл - в этом файле мы и искали общие мотивы.}

\par{Далее проверку на совпадения вышестоящих областей мы засунули все вышестоящие области хоста и его фагов в программу
MEME Suite. Как и ожидалось, мы получили множество профилей, которые были общими у фагов и их хоста. }

\par{Также MEME-Suite для каждого профиля мотива дала нам PFM (Position Frequency Matrix), поэтому у нас была возможность
сравнить сами мотивы между собой. У нас появилась гипотеза, что мотивы фагов, атакующих одного хоста, могут быть также
схожи и как-то ировать между собой. Мы решили посмотреть корреляцию 2D массивов - потому что для каждой позиции в
матрице частот у нас было 4 значения - на 4 возможных нуклеотида.}

\newpage
\begin{center}
    \item \section{Результаты и Обсуждение}
    \item \subsection{MEME-Suite Output}
\end{center}

\begin{itemize}
    \item Логотипы мотивов - Логотипы последовательностей представляют собой графическое представление информационного
    содержимого, хранящегося в системе множественного выравнивания последовательностей (MSA), и обеспечивают
    компактное и интуитивно понятное представление специфичного для положения нуклеотидного состава связывающих
    мотивов, активных сайтов и т.д. в биологических последовательностях. (Reference: Thomsen, M.C., \& Nielsen, M.
    2012. Nucleic Acids Res. 40(Web Server issue):W281-287).
   
    \item Матрица частот нуклеотидов - считается путем подсчета вхождений каждого нуклеотида в каждой позиции. Один мотив встречается в последовательностях несколько раз и с наперёд заданной точностью мы считаем, что это подходящий для нас мотив.
    
    \item PSSM (Position-Specific Scoring Matrices) - The position-specific scoring matrix corresponding to the motif
    is printed for use by database search programs such as MAST. This matrix is a log-odds matrix calculated by taking
    100 times the log (base 2) of the ratio p/f at each position in the motif where p is the probability of a
    particular letter at that position in the motif, and f is the background frequency of the letter (given in the
    command line summary section.) This is the same matrix that is used above in computing the p-values of the
    occurrences of the motif in the Occurrences of the Motif and Block Diagrams of Motif Occurrences sections. The
    scoring matrix is printed "sideways"\-\-columns correspond to the letters in the alphabet (in the same order as
    shown in the simplified motif) and rows corresponding to the positions of the motif, position one first. The
    scoring matrix is preceded by a line starting with "log-odds matrix:" and containing the length of the alphabet,
    width of the motif, number of characters in the training set, the scoring threshold (obsolete) and the motif
    E-value. \cite{memeres} 
    
    \item host\_phage\_results\_total.tsv - таблица, в которой содержатся данные о мотивах: последовательность, длина, число сайтов, ic, llr, e\_value, название файла с PSSM, название файла с PFM, регулярное выражения для мотива и причина остановки.
\end{itemize}


\begin{center}
\item \subsection{Обсуждения}
\end{center}

    \par{Прежде всего, мы уже знаем, что у бактерии и у хоста повторяются -35 и -10 промоторные элементы. Нам они не
    были интересны. Мы нашли возможные последовательности для -35 и для -10 участков \cite{-35,-10,-35-10,-35-10wiki}.
    Нашли и удалили их из нашего датасета с помощью регулярных выражений \cite{re}. }
    
    \par{Нам было интересно положение мотивов относительно начала транскрипции. Как уже было сказано, мы выделили 50 
    нуклеотидов - перед стартом транскрипции. Опять же, с помощью регулярных выражений, мы нашли положения мотивов в
    каждом 50-нуклеотидном регионе, в котором он встречается. Для каждого мотива мы посчитали среднее значение для его
    левого конца относительно начала транскрипции, для его правого конца и среднеквадратичную ошибку для обоих
    концов. Результаты записали в файл $ host\_phage\_results\_total\_updated.tsv $}
    
    \par{Дальше для каждого мотива мы нашли где контретно он находится и добавили в host\_phage\_results\_total.tsv
    отдельную колонку "location". В этой колонке присутствуют значения BOTH - если мотив был найден и у бактерии и у
    фага, либо "PHAGE", если мотив был найден только и фага или "HOST", если мотив был найден только у бактерии. К
    нашему удивлению некоторые мотивы не были найдены нигде - для таких мы в этой колонке выписали "NOT FOUND".}
    
    \par{Нам также стало интересно, коррелируют ли как-то между собой мотивы бактериофагов, которые атакуют одну бактерию, а также есть ли корреляция между мотивами бакториофагов из одной таксономической категории. Для начала мы хотели посчитать корреляцию Пирсона между PFM этих мотивов. Но PFM - двумерный массив, мы не обнаружили в известных библиотеках Python подходящей функции для подсчётка корреляции между 2D-массивами. Мы написали собственный скрипт для предварительного подсчёта такой корреляции $preliminary_correlation.py$, используя трюки отсюда \cite{stackoverflow}. Результаты на данном этапе трудно интерпретировать, но можно заметить, что данные и вправду как-то коррелируют.}
    
    \par{Также, как я уже написал выше, мы удивлись, когда обнаружили, что некоторые мотивы не были найдены ни в каком из представленных вышестоящих регуляторных регионах. Но мы искали сами мотивы внутри них. Также во многих мотивах можно заметить большое количество букв W, N, Y и ещё некоторых, которые не входят в состав привычных нам A,T,G и C для нуклеотидов. Дело в том, что на одном месте в мотиве в разных последовательностях могут стоять разные нуклеотиды. И для таких двойственных мест MEME-Suite использует специальный алфавит, в котором одна буква заменяет все возможные основания на данной позиции \cite{memealphabet}. }
    
    \par{Но когда мы начали анализировать эти последовательности, мы заметили, что некоторые регулярные выражения и двойственные буквы, которые должны им соответствовать, не совпадают с тем, что дано в алфавите \cite{memealphabet}. Мы решили протестировать нашу выборку в ещё одно программе по поиску собственных мотивов. К сожаления, пока что мы не нашли подходящей программы.
    \begin{itemize}
        \item sarus-2.0.2 ищет только известный мотив в нескольких последовательностях (а нам нужно в них самих находить новые). 
        \item Weeder подходит под наши критерии, но он ищет только мотивы длиной 6, 8 и 10 нуклеотидов - для нас это маловато. мы же искали мотивы длиной от 10 до 50 нуклеотидов.
        \item Phyloscan, motifmatchr и много других пакетов для R всё ещё мне недоступны. Я с R впервые познакомился в этом семестре, поэтому я не успел пока к нему приспособиться, чтоби использовать в дипломе какие-то пакеты для R.
        \item MotifScan и много подобных ему ищут лишь мотивы из известных баз данных. Здесь мы столкнулись с той же проблемой, что и при использовании sarus-2.0.2
    \end{itemize}
    }
    
    \par{Кратко описал результаты. Видно сейчас, что работа на будущее у нас есть: мы хотим написать свой скрипт для корелляционого анализа 2D массивов и + проверить наши мотивы в другой прогаммы для уверенности. А дальше как пойдёт}
    
\newpage
\begin{center}
\item \section{Выводы}
\end{center}

\newpage
\begin{thebibliography}{9}
    \addcontentsline{toc}{section}{\refname}
    \bibitem{hinton1} Stitt B, Hinton DM: Regulation of middle-mode transcription.In Molecular biology  of  bacteriophage  T4.Edited by: Karam JD, Drake J, Kreuzer KN, Mosig G, Hall D, Eiserling F, Black L, Spicer E, Kutter E, Carlson K, Miller ES.Washington, D.C.: American Society for Microbiology; 1994:142-160
    \bibitem{hinton2} Brody E, Rabussay D, Hall D:Regulation of transcription of prereplicative genes. In Bacteriophage  T4.Edited by: Mathews CK, Kutter EM, Mosig G,Berget PB. Washington, D. C.: American Society for Microbiology;1983:174-183
    \bibitem{hinton3} Hinton DM, Pande S, Wais N, Johnson XB, Vuthoori M, Makela A, Hook-Barnard I: Transcriptional takeover by sigma appropriation: remodelling of the sigma70 subunit of Escherichia coli RNA polymerase by the bacteriophage T4 activator MotA and co-activator AsiA.Microbiology2005,151:1729-1740.
    \bibitem{hinton4} Miller ES, Kutter E, Mosig G, Arisaka F, Kunisawa T, Ruger W: BacteriophageT4 genome. Microbiol  Mol  Biol  Rev2003,67:86-156
    \bibitem{hinton5} Wilkens K, Ruger W: Transcription from early promoters. In Molecular Biology  of  Bacteriophage  T4.Edited by: Karam JD, Drake JW, Kreuzer KN,Mosig G, Hall DH, Eiserling FA, Black LW, Spicer EK, Kutter E, Carlson K,Miller ES. Washington, D. C.: American Society for Microbiology;1994:132-141.
    \bibitem{hinton6} Weisberg R, Hinton DM, Adhya S: Transcriptional Regulation in Bacteriophage. In Encyclopedia  of  Virology.3 edition. Edited by: Mahy BWJ,van Regenmortel MHV. Oxford: Elsevier; 2008:174-186, 174-186.
    \bibitem{hinton7} Hinton DM: Transcription from a bacteriophage T4 middle promoter using T4 motA protein and phage-modified RNA polymerase.J  BiolChem 1991, 266:18034-18044
    \bibitem{hinton8} Schmidt RP, Kreuzer KN: Purified MotA protein binds the -30 region of a bacteriophage T4 middle-mode promoter and activates transcription invitro. J  Biol  Chem1992, 267:11399-11407
    \bibitem{hinton9} Stevens A: New small polypeptides associated with DNA-dependent RNA polymerase of Escherichia coli after infection with bacteriophage T4. Proc  Natl  Acad  Sci  USA1972, 69:603-607.
    \bibitem{hinton10} Stevens A:An inhibitor of host sigma-stimulated core enzyme activitythat purifies with DNA-dependent RNA polymerase of E. coli followingT4 phage infection.Biochem  Biophys  Res  Commun1973,54:488-493.
    \bibitem{hinton11} Ouhammouch M, Orsini G, Brody EN: The asiA gene product of bacteriophage T4 is required for middle mode RNA synthesis.J  Bacteriol1994, 176:3956-3965
    \bibitem{hinton12} Ouhammouch M, Adelman K, Harvey SR, Orsini G, Brody EN: BacteriophageT4 MotA and AsiA proteins suffice to direct Escherichia coli RNA polymerase to initiate transcription at T4 middle promoters. Proc  NatlAcad  Sci  USA1995, 92:1451-1455.
    \bibitem{hinton13} Hinton DM, March-Amegadzie R, Gerber JS, Sharma M: Bacteriophage T4middle transcription system: T4-modified RNA polymerase; AsiA, a sigma70 binding protein; and transcriptional activator MotA.Methods  Enzymol1996, 274:43-57
    \bibitem{virus-host} https://www.genome.jp/virushostdb/ 
    \bibitem{biopython} https://biopython.org/
    \bibitem{python} https://docs.python.org/3/
    \bibitem{jupyter} https://jupyter.org/
    \bibitem{github} https://github.com/poolsar42/phages-and-hosts
    \bibitem{entrez} https://www.ncbi.nlm.nih.gov/books/NBK25501/
    \bibitem{hinton} Hinton D.M. Transcriptional control in the prereplicative phase of T4 development // Virology
    journal. 2010. V. 7. P. 289. https://doi.org/10.1186/1743-422X-7-289
    \bibitem{puxty-evanx} Puxty, R.J., Evans, D.J., Millard, A.D. et al. Energy limitation of cyanophage development:
    implications for marine carbon cycling // ISME J. 2018, 12, 1273–1286. https://doi.org/10.1038/s41396-017-0043-3
    \bibitem{lomsad} Lomsadze, A., Gemayel, K., Tang, S., and Borodovsky, M. Modeling leaderless transcription and
    atypical genes results in more accurate gene prediction in prokaryotes // Genome research,  2018. 28(7), 1079–1089.
    https://doi.org/10.1101/gr.230615.117
    \bibitem{bailey} Bailey, T.L., Johnson, J., Grant, C.E., and Noble, W.S. The MEME Suite // Nucleic acids research,
    2015, 43(W1), W39–W49. https://doi.org/10.1093/nar/gkv416
    \bibitem{ryasik} Ryasik, A., Orlov, M., Zykova, E., Ermak, T., Sorokin, A. Bacterial promoter prediction: Selection
    of dynamic and static physical properties of DNA for reliable sequence classification // 2018
    https://www.worldscientific.com/doi/abs/10.1142/S0219720018400036
    \bibitem{genomics} https://www.sciencedirect.com/science/article/pii/S1369527403001152?via\%3Dihub
    \bibitem{structure} https://www.pnas.org/content/108/28/E288.full
    \bibitem{networks} https://www.sciencedirect.com/science/article/pii/S0966842X12002004
    \bibitem{memeres} https://kodomo.fbb.msu.ru/~partyhard/term4/pr9/meme\_out\_oops/meme.html\#pssm1
    \bibitem{-35} https://www.addgene.org/mol-bio-reference/promoters/
    \bibitem{-35-10} https://www.sciencedirect.com/topics/biochemistry-genetics-and-molecular-biology/pribnow-box
    \bibitem{-10} https://www.cell.com/cell/fulltext/S0092-8674(11)01339-0
    \bibitem{-35-10wiki} https://en.wikipedia.org/wiki/Promoter\_(genetics)\#Bacterial
    \bibitem{re} https://docs.python.org/3/library/re.html\#module-re
    \bibitem{stackoverflow} https://stackoverflow.com/questions/30143417/computing-the-correlation-coefficient-between-two-multi-dimensional-arrays
    \bibitem{memealphabet} https://meme-suite.org/meme/doc/alphabet-format.html\#standard\_DNA
\end{thebibliography}

\end{document}
